\section{Normalization and finite identities}\label{sec:preliminaries}

Write $D(c,r)=\{z\in\C:|z-c|<r\}$ for the open disk, and
$t_+=\max(t,0)$ for real $t$.
For vectors, $\|\cdot\|$ is the Euclidean norm; for matrices it is
the associated operator norm. We use $A^*$ for the conjugate transpose,
$\|A\|_F^2=\sum_{i,j}|A_{ij}|^2$ for the squared Frobenius norm,
and $A\ge0$ for a Hermitian positive semidefinite matrix.
The symbols $I$ and $\mathbf1$ denote the identity matrix and the
all-ones vector of the required size; $\tr$ and $\diag$ denote
the trace and diagonal-matrix operations.

Write $N=n-1$, and list the zeros of $p$ as $z_1,\ldots,z_n$ and
the critical points as $\xi_1,\ldots,\xi_N$. A \emph{strict hole}
is a point $a\in\conv Z(p)$ for which $|a-\xi_j|>1$ for every $j$.
The negation of Theorem~\ref{thm:main} is precisely the existence of
such a point. We may multiply $p$ by a nonzero constant whenever
convenient.

\begin{lemma}[Gauss--Lucas]\label{pr:gauss-lucas}
Every critical point belongs to $\conv Z(p)$.
\end{lemma}
\begin{proof}
If $w\notin\conv Z(p)$, a separating line and a rotation give
$\Re(e^{-i\theta}(w-z_k))>0$ for every $k$. Thus every
$e^{i\theta}/(w-z_k)$ has positive real part. Their sum is
$e^{i\theta}p'(w)/p(w)$ and cannot vanish. A zero shared by $p$
and $p'$ already lies in the zero hull.
\end{proof}

\begin{lemma}\label{pr:edge-reduction}
Suppose $n\ge6251$ and all zeros lie in $\overline\D$. If a strict
hole exists, after a rotation there is a strict hole $a$ with
$0<a<1$ on an actual supporting edge of a two-dimensional zero hull.
For every critical point,
\begin{equation}\label{pr:critical-halfplane}
 \Re\xi_j<a/2.
\end{equation}
There is an actual simple rightmost zero $b=x+iy$ with $x\ge a$.
After conjugation we may take $y\ge0$.
\end{lemma}
\begin{proof}
Gauss--Lucas implies $|\xi_j|\le1$, so zero cannot be a strict
hole. Rotate a strict hole onto the positive real axis. Then
\[
 2a\Re\xi_j<a^2+|\xi_j|^2-1\le a^2.
\]
Let $A$ be the last point of the zero hull on that ray. For
$a\le A\le1$,
\[
 |A-\xi_j|^2-|a-\xi_j|^2
 =(A-a)(A+a-2\Re\xi_j)\ge0.
\]
Thus the hole persists at $A$. If $A$ is an actual zero, it is either
multiple, in which case it is critical, or simple. In the latter
case Proposition~\ref{ar:energy} gives
$\sum_j|A-\xi_j|^{-2}\ge N$, contrary to all the distances being
greater than one.

A hull consisting of one point has no hole. If the hull is a segment,
Gauss--Lucas puts all critical points on it. The distance between any
two consecutive distinct critical points is at most two, since the
whole segment lies in the unit disk. Their closed unit intervals
therefore have connected union. Each endpoint of the hull belongs
to that union by the preceding actual-zero argument, so the whole
segment is covered. In the remaining case $A$ lies on a supporting
edge whose endpoints are actual zeros. Strict convexity of the unit
disk excludes $A=1$. Rename $A$ as $a$ and repeat the displayed
half-plane inequality at this point. A zero maximizing real part
has real part at least $a$; it cannot be multiple by
\eqref{pr:critical-halfplane}. Conjugation makes its imaginary part
nonnegative.
\end{proof}

For an actual simple rightmost zero $b$, put
\begin{equation}\label{pr:field-notation}
 v_j=(b-\xi_j)^{-1},\quad S=\sum_jv_j=\frac{p''(b)}{p'(b)},
 \quad H=|S|,\quad Q=\sum_j|v_j|^2,\quad
 d=\min_j|b-\xi_j|.
\end{equation}
Rightmostness and Gauss--Lucas give $\Re v_j\ge0$ and
$|v_j|\le1/d$. The equality $S=p''(b)/p'(b)$ counts the complete
critical divisor. If the other original zeros are $u_1,\ldots,u_N$,
factoring $p(z)=(z-b)g(z)$ gives the further identity
\begin{equation}\label{pr:actual-root-identity}
 S=2\sum_{k=1}^N\frac1{b-u_k}.
\end{equation}

\begin{lemma}\label{pr:radial}
The data in \eqref{pr:field-notation} satisfy
\begin{equation}\label{pr:radial-identity}
 N=2\Re(bS)+(1-|b|^2)Q
       -\sum_j\frac{1-|\xi_j|^2}{|b-\xi_j|^2}.
\end{equation}
If all zeros lie in $\overline\D$, $a>0$ is a strict hole, and
$b=x+iy$ is rightmost with $x\ge a$,
then $\Re S>0$ and
\begin{equation}\label{pr:hole-reciprocal}
 Q<\frac{2\Re S}{2x-a}.
\end{equation}
\end{lemma}
\begin{proof}
For each $\xi$ expand
\[
 2\Re\frac b{b-\xi}
 =1+\frac{|b|^2-|\xi|^2}{|b-\xi|^2}
\]
and sum. For the second assertion, Gauss--Lucas and the strict hole give
$2a\Re\xi<a^2+|\xi|^2-1\le a^2$. Thus
$\Re(b-\xi)^{-1}=(x-\Re\xi)/|b-\xi|^2
 >(x-a/2)/|b-\xi|^2$. Summing proves the assertion.
\end{proof}

\begin{lemma}\label{pr:primitive}
In the disk $|z|<d$, with the branch of each logarithm equal to zero
at the origin,
\[
 \frac{p'(b+z)}{p'(b)}=\prod_{j=1}^N(1+zv_j),\qquad
 p(b+z)=p'(b)\int_0^z\prod_{j=1}^N(1+tv_j)\,dt.
\]
For every integer $r\ge2$ and $|z|<d$,
\begin{equation}\label{pr:complete-log-tail}
 \left|\sum_j\Log(1+zv_j)
 -\sum_{k=1}^{r-1}\frac{(-1)^{k+1}z^k}{k}\sum_jv_j^k\right|
 \le\frac{Q|z|^r}{r d^{r-2}(1-|z|/d)}.
\end{equation}
\end{lemma}
\begin{proof}
The product is the factorization of $p'$; integration uses $p(b)=0$.
The logarithms are analytic because $|zv_j|<1$. For $k\ge r$,
$\sum_j|v_j|^k\le Qd^{2-k}$ and $1/k\le1/r$.
Summing the entire remaining geometric series proves the bound.
\end{proof}

When $H>0$ it is useful to write
\begin{equation}\label{pr:normalized-moments}
 D=Hd,\qquad \mu=D^{-1},\qquad q=Q/H^2,\qquad
 \sigma=S/H,\qquad w_j=v_j/H.
\end{equation}
Then $|\sigma|=1$, $\sum w_j=\sigma$,
$\sum|w_j|^2=q$, and $\max|w_j|=\mu$. For $r>0$, the
polynomial $p_r(z)=p(rz)$ at the corresponding zero $b/r$ has
field data $d/r,rS,rH,r^2Q$. Thus $D,q,\mu,\sigma$ are invariant
under positive real rescaling. Rescaled coordinates will be specified
when used.
